\documentclass[11pt]{article}
\usepackage{icmlko}
\begin{document}

\title{성적표가 불공정했다 --- 성과가 벌어지는 폭으로 표준화하니 만화가 아홉째에서 둘째로 올라간다}
\author{월드모델 실험실}
\date{노트 534 · 2026년 8월 4일}
\maketitle

\begin{abstract}
노트 533 이 도메인마다 \textbf{성과가 벌어지는 폭}이 크게 다름을 봤다
(후보 다섯 중 $1$--$5$ 등 배수: 모바일 $296$ 배 $\cdot$ 게임 $117$ 배 대
\textbf{만화 $3.5$ 배 $\cdot$ 세계애니 $3.7$ 배}). 폭이 좁으면 어떤 모델도
가르기 어렵다 --- \textbf{그러면 도메인별 정확도를 나란히 놓는 것 자체가
불공정하다.}

\begin{center}
\begin{tabular}{lrrrr}
\hline
도메인 & 성과 폭 & 날 정확도 & 표준화 & 차 \\
\hline
\textbf{만화} & $4.4$ 배 & $61.5\%$ (9위) & $\mathbf{74.7\%}$ (2위) & $\mathbf{+13.2}$ \%p \\
\textbf{세계애니} & $4.8$ 배 & $70.4\%$ (4위) & $\mathbf{81.6\%}$ (1위) & $+11.1$ \%p \\
모바일 & $651$ 배 & $72.5\%$ (2위) & $66.0\%$ (6위) & $-6.5$ \%p \\
게임 & $240$ 배 & $72.9\%$ (1위) & $68.2\%$ (4위) & $-4.7$ \%p \\
시장팝업 & $14$ 배 & $54.2\%$ (11위) & $54.5\%$ (11위) & $+0.3$ \%p \\
\hline
\end{tabular}
\end{center}

사전 등록 둘 다 통과 --- 성과 폭 $\leftrightarrow$ 날 정확도 순위상관
$+0.691$, 날 순위 $\leftrightarrow$ 표준화 순위 $+0.609$,
그리고 \textbf{만화 $+7$ 칸 $\cdot$ 세계애니 $+3$ 칸}.
\end{abstract}

\begin{figure}[t]
\centering
\includegraphics[width=\textwidth]{figs/f1.pdf}
\caption{\textbf{가} --- 성과 폭과 날 정확도. \textbf{나} --- 순위 이동.
\textbf{다} --- 외삽 검사.}
\end{figure}

\section{표준화는 연령표준화와 같다}

전 도메인을 합친 \emph{배수 구간별 쌍 비중}을 표준 인구로 잡고, 도메인마다
그 구간 안 정확도를 그 비중으로 다시 가중한다.

$$\text{표준화 정확도}_d = \sum_b w^{\text{표준}}_b \cdot a_{d,b}$$

폭이 넓은 도메인(모바일 $\cdot$ 게임)은 쉬운 큰 배수 쌍을 많이 갖고 있어
날 정확도가 부풀고, 폭이 좁은 도메인(만화 $\cdot$ 세계애니)은 그 반대다.

\section{외삽이 아니다 --- 전혀 다른 방법으로 검산했다}

의심할 것이 있었다. 만화의 $5$ 배 이상 칸은 자기 쌍의 $6.1\%$ 뿐인데 표준
무게는 $39\%$ 다 --- \textbf{얇은 칸에 기대 외삽하는 것 아닌가?}

\begin{quote}
\textbf{쌍 수로 보면 $2{,}031$ 쌍이다}(세계애니 $3{,}293$). 비중이 낮을 뿐
표본은 넉넉하다.
\end{quote}

그리고 \textbf{외삽이 $0$ 인 비교}를 따로 했다 --- 모든 도메인이 두꺼운 칸
($2\!\sim\!5$ 배, 최소 $310$ 쌍)에서만 견주기.

\begin{center}
\begin{tabular}{lrr}
\hline
도메인 & $2\!\sim\!5$ 배 구간 정확도 & 쌍 수 \\
\hline
세계애니 & $\mathbf{84.5\%}$ & $13{,}740$ \\
만화 & $\mathbf{66.8\%}$ & $10{,}300$ \\
게임 & $64.4\%$ & $3{,}688$ \\
애니 & $64.3\%$ & $51{,}471$ \\
웹툰 & $64.0\%$ & $82{,}157$ \\
모바일 & $60.8\%$ & $19{,}037$ \\
시장팝업 & $54.1\%$ & $2{,}704$ \\
\hline
\end{tabular}
\end{center}

\begin{quote}
\textbf{순위상관 --- 날 정확도 $\leftrightarrow$ 고정 구간 $+0.527$,
표준화 $\leftrightarrow$ 고정 구간 $\mathbf{+0.955}$.} 표준화는 외삽이
아니라 \emph{같은 것}을 재고 있었다.
\end{quote}

\section{그래서 무엇이 뒤집히나}

\begin{quote}
\textbf{모바일은 날 정확도 $2$ 위인데 같은 배수 구간에서는 $8$ 위다.}
``모바일이 잘 맞는다''는 \textbf{성과 폭이 $651$ 배여서}였다.
반대로 \textbf{만화는 우리가 제일 못 맞히는 시장이 아니라 둘째로 잘 맞히는
시장}이다 --- 성과가 $4.4$ 배 안에 몰려 있는데도 $66.8\%$ 를 낸다.
\end{quote}

이것은 노트 505 와 겹쳐 읽어야 한다. 거기서 만화는 ``자기 천장의 $2.00$ 배
--- 제일 많이 받고 유일하게 안 준다''였다. 두 관측이 모순이 아니다 ---
\textbf{만화는 남에게서 많이 받아서 자기 폭 안에서 잘 맞히고, 그 대신 남에게
줄 것은 없다.}

\textbf{시장팝업만 어느 셈법에서도 꼴찌다}(날 $54.2$ $\cdot$ 표준화 $54.5$
$\cdot$ 고정 구간 $54.1$). 노트 530 의 태거 진단(유보 $126$ 행이 $27$ 가지
설명으로 무너진다)이 여기서 한 겹 더 단단해진다 --- \textbf{시장팝업의 낮은
점수는 성과 폭 탓이 아니다.}

\section*{조항}

\begin{quote}
\textbf{도메인별 성적표는 성과 폭으로 표준화해서 읽는다.} 날 정확도를 나란히
놓는 것은 시험 난이도가 다른 반의 평균 점수를 견주는 것과 같다. 표준화가
얇은 칸에 기대는지 의심되면 \textbf{모두가 두꺼운 칸 하나에서만} 견주는
검산을 함께 낸다.
\end{quote}

\section*{남은 것}

\begin{itemize}
\item \textbf{노트 508 의 수집 우선순위를 다시 본다.} 거기서 모바일이 $1$ 순위
(행당 값어치 $-0.0610$)였는데, 그것은 \emph{판 기여}라 이 노트와 다른 축이다.
다만 \textbf{모바일이 실제로는 잘 맞히는 시장이 아니라는 것}은 함께 읽어야 한다.
\item \textbf{세계애니 $84.5\%$ 는 어디서 오나} --- $2\!\sim\!5$ 배 구간에서
전 도메인 최고다. 무엇이 그 시장을 쉽게 만드는지 아직 모른다.
\item \textbf{판($\rho$)도 폭에 흔들리나} --- 이 노트는 쌍 정확도만 표준화했다.
순위 규약 자체가 같은 편의를 갖는지는 안 쟀다.
\end{itemize}

\end{document}
